\clearpage\thispagestyle{empty}
\pdfbookmark[0]{Sources for Gauss’s Letters on De Vico’s Comet (October and November 1844)}{cometdevico-sources}
\begin{center}\textbf{Sources for Gauss’s Letters on De Vico’s Comet (October and November 1844)}\end{center}
\noindent This AI translation follows the extracts of 15 October and 11 November 1844 on pp.~469--470 of \href{https://gdz.sub.uni-goettingen.de/id/PPN236007467}{\textit{Werke}, vol.~VI} (Göttingen, 1874). Carl Friedrich Gauss wrote to Heinrich Christian Schumacher. In \textit{Astronomische Nachrichten}, vol.~XXII, no.~521, \href{https://api.digitale-sammlungen.de/iiif/image/v2/bsb10049447_00147/full/1800,/0/default.jpg}{cols.~277--278}, the journal signs the extracts \textit{Gauss} and identifies the comet as De Vico’s in the ephemeris heading. The issue is dated 28 November 1844 on \href{https://api.digitale-sammlungen.de/iiif/image/v2/bsb10049447_00148/full/1800,/0/default.jpg}{col.~280}.
\par\smallskip\noindent Goldschmidt calculated the elliptic elements. The collected text ends with the sentence introducing the ephemeris. The journal prints the ephemeris on cols.~277--280, with a closing explanation of its brightness column. Translation and typesetting produced by AI.
\clearpage
